\subsection{Experimental protocol}
\label{sec:benchmarks}

\subsubsection{Optimization sequences}
The body-fitted sequences contain 16 distinct targets. Two moving
Gaussian profiles define each desired state, and a pulse changes its
amplitude during a transient trajectory.
The six primary cases comprise steady and four-slab transient
transformer problems, together with the corresponding bore-in-block problems at
two nested refinement levels. They use rank 100 and five complete
repetitions per solver. Section~\ref{sec:v20_results} adds larger meshes,
longer trajectories and changing operating conditions under the procedure of
Section~\ref{sec:final_procedure}.

The operating sequences use one reference built at nominal conditions. The
transformer day chains 24 hourly eight-slab windows through a residential load
cycle, with the winding source proportional to the squared per-unit load. The
overload sequence applies four emergency load profiles to the eight-slab transformer: a
step to 1.3 per unit, a ramp to 1.5 per unit, a 1.8 per-unit pulse and an alternation
between 0.8 and 1.3 per unit. Each profile is paired with one target and four bounds
from 0.7 to 0.85, which gives 16 queries. The drive
cycle chains 16 four-slab windows on Bore~2, with the interface heat flux and swirl
speed changing between phases. Each chained window starts from the final state of
the preceding one. The off-design sequences contain 16 queries, four targets at each
of the velocity scales 0.5, 0.75, 1.5 and 2, and reassemble the operator for every
query. The shuffled sequences permute the query order.

Each Cartesian sequence contains 16 desired states indexed by
$\theta_m=m/15$, $m=0,\ldots,15$. Four Gaussian profiles have
unequal weights and query-dependent centers. Rotation and vertical
translation change their locations across queries. Transient targets
also rotate and pulse over physical time. Hashes verify that the targets are
distinct and identical between methods and repetitions.

The Cartesian steady bound is approximately $\psi=1.950\times10^{-4}$,
and its transient bound is approximately $\psi=2.685\times10^{-4}$.
The configurations retain the full-precision values.
A direct-PDAS calibration selected each bound once; the
resulting value stays fixed under spatial refinement and changes in
the number of time slabs. Uniform steps partition the fixed horizon
$T_f=0.1$ in the performance study. Unequal steps and nonzero
initial data appear in the independent discretization checks.

\subsubsection{Four complete-optimization solvers}
The first method is ordinary Jacobi-preconditioned CG with an
efficient rank-zero path. The second is RefDef with
\eqref{eq:reference_restriction}. Both use the same operator
application, precision and residual verification. The third uses
the same deflated-CG kernel with history-based recycling. It starts
with empty history. After an inner solve meets the residual criterion, its candidate
pool contains the deployed coarse vectors and the last $w$ new
CG search directions, with $w=r$. The directions are copied during
the solve, and the space remains fixed until that solve ends.

For physical candidate matrix $V$ and current Jacobi diagonal $D$,
recycling forms an orthonormal basis $\widehat V$ of $D^{1/2}V$.
It selects the lowest Ritz vectors of
$\widehat V^\top D^{-1/2}BD^{-1/2}\widehat V$ and maps them back
by $D^{-1/2}$. QR and a small-factor SVD use a relative rank threshold
of $10^{-12}$. At most $r$ physical vectors are retained.
Zero extension transfers them to the next inactive set.
Capture, selection, transfer and cached products all enter the
cost and memory measurements. A failed target clears the history.

The fourth method is AmgX-preconditioned CG with classical AMG. Configuration and
resource handles persist throughout a sequence, and each inactive matrix receives
new matrix and vector handles, a solver and an AMG hierarchy.

\subsubsection{Warm starts and accuracy}
The primary policy warm-starts the outer and inner solves.
The first query starts with all state DOFs active.
Subsequent queries receive the preceding solution and active
set after the KKT check. Within PDAS, each verified inactive solution updates
the full state used by the next inner solve. Every method has equivalent
access to this solution history. The cold-start variant resets the
active set to all DOFs for each query and uses zero
inner initial guesses. Each complete repetition starts without
solution or recycling history.

All computations use double precision. The final inactive-system
criterion is an independently recomputed relative residual of
$10^{-10}$. CG triggers candidate verification at a relative
recurrence residual of $10^{-11}$; AmgX uses a native absolute
target of $10^{-11}\|b\|_2$. For a zero right-hand side, the
corresponding absolute targets apply. The common independent
check determines acceptance. Each Cartesian inner solve has a cap of
10000 iterations. The body-fitted study uses a cap of 50000 iterations.
PDAS has a cap of 100 steps in both studies, raised to 400 for Bore~4.

Set $g=MY-\ell$ and $s=\max(1,\|\ell\|_\infty)$.
PDAS requires each of the following quantities to be at most
$10^{-8}$:
\begin{equation}\label{eq:cht_kkt_acceptance}
\begin{gathered}
 \|g+\lambda\|_\infty/s,\qquad
 \|[Y-\psi]_+\|_\infty,\qquad
 \|[-\lambda]_+\|_\infty/s,\\
 \|\lambda\odot(\psi-Y)\|_\infty/s,\qquad
 \|Y-\min(\psi,Y-g)\|_\infty/s.
\end{gathered}
\end{equation}
The primal-feasibility quantity is absolute. The other four use
the stated common normalization. Transient runs additionally
check \eqref{eq:control_adjoint} through independent slabwise
forward and transpose actions. Caps, cycles, memory limits,
exceptions and residual failures are recorded separately. The reported
times contain only sequences that satisfy these criteria at every target.

\subsubsection{Time, memory and reproducibility}
Each method and repetition runs in a fresh process on one GPU
with four central processing unit (CPU) threads. A common small warmup initializes both
GPU libraries. The method order rotates between repetitions.
The complete sequence interval starts before operator assembly
and reference or resource construction. It includes every target,
all PDAS updates and inactive solves, matrix and basis transfers,
verification and cleanup. Disjoint timing components sum to this
measured total. Reference construction is charged once, to the first
query of each sequence.

Calibration, process initialization, warmup and termination are recorded
separately; the body-fitted models use fixed bounds and need no calibration. A
separate process samples host and GPU memory.

Each case compares RefDef with the fastest tested GPU iterative
alternative: among Jacobi-CG, recycling and AmgX, the method with the lowest
median complete time over its converged repetitions, provided at least one
repetition converges at every target. We call it the fastest alternative, and the
speedup divides its median time by that of RefDef. The main-text
comparisons are those of Tables~\ref{tab:mesh_complete} and
\ref{tab:v20_scale}--\ref{tab:v20_operation}. The direct and block-in-time CPU
solvers are compared separately; the block-in-time solver applies CG with block-Jacobi
sweeps over the slab blocks of the trajectory Hessian. Most comparisons run all methods
and repetitions within one job on one GPU; the table captions identify comparisons
assembled from separate jobs. Repeated measurements of one Bore~3 configuration in
separate jobs show between-job timing variation. The reported envelopes describe the
sampled repetitions, so small speedup margins should be interpreted within that
measurement scope.

Source identifiers, configurations, target and matrix hashes, compressed
active sets and scripts that regenerate every table and figure accompany the
timing records.

\subsubsection{Timing procedures}
\label{sec:final_procedure}
Two procedures share Algorithm~\ref{alg:revision_deflation} and the accuracy
criteria above (Table~\ref{tab:procedures}). The
per-update assembly procedure forms each inactive Hessian by sparse products after
every active-set update; the primary comparisons of Section~\ref{sec:mesh_results}
use it, as the original body-fitted implementation does. The sliced-Hessian procedure
of Algorithm~\ref{alg:final_procedure} assembles the full reduced Hessian once for
each operator and extracts every inactive principal submatrix from it; the comparisons
of Section~\ref{sec:v20_results} use it. The extraction removes the per-update sparse
products from the setup of every method.
Reference meshes, ranks and temporal coverage are prescribed settings of each
comparison. One reference, computed before the first query at the nominal flow, is
used for every query of a sequence.

Targets, loads, bounds and chained initial states change only the right-hand side
or the constraints and leave the full reduced Hessian unchanged; each active-set update
extracts its current inactive principal submatrix. A change of flow scale reassembles
the full reduced Hessian. The comparisons use Algorithm~\ref{alg:revision_deflation}
with the coarse correction and residual replacement. The \href{https://github.com/teeratornk/deflation-example/blob/ba682dd05b168d2f43ff0c58050813deee785223/docs/mesh-showcases.md#choose-the-numerical-procedure}
{reproduction examples} give the commands for each procedure.

\begin{table}[pos=htbp]
\centering\footnotesize\setlength{\tabcolsep}{3pt}
\caption{Timing procedures and their prescribed reference settings. Both procedures
solve each inactive system with Algorithm~\ref{alg:revision_deflation}, including its
residual replacement, and restrict the reference on the GPU. Complete temporal coverage
keeps every temporal factor of each retained spatial mode; truncated coverage keeps the
lowest spatial--temporal pairs up to the total rank. Temporal coverage does not apply to
steady cases.}
\label{tab:procedures}
\begin{tabularx}{\textwidth}{@{}>{\raggedright\arraybackslash}p{0.15\textwidth}>{\raggedright\arraybackslash}p{0.12\textwidth}>{\raggedright\arraybackslash}p{0.15\textwidth}>{\raggedright\arraybackslash}X>{\raggedright\arraybackslash}p{0.22\textwidth}@{}}
\toprule
Procedure & Comparisons & Inactive Hessian & Reference mesh and rank & Temporal coverage \\
\midrule
Per-update assembly & Table~\ref{tab:mesh_complete} & Sparse products at every update &
Packaged mesh, total rank 100 & Truncated with four slabs; the steady cases have no
temporal factors \\
Sliced Hessian & Table~\ref{tab:v20_scale} (steady) & Principal submatrix of the
preassembled Hessian & Closed form on Cartesian grids, rank 100; packaged mesh for the
transformer and Bore~1, rank 100; first refinement for Bore~2 and Bore~3 and second
refinement for Bore~4, rank 200 & No temporal factors \\
Sliced Hessian & Table~\ref{tab:v20_transient} (trajectories) & Principal submatrix of
the preassembled Hessian & Closed form on Cartesian grids; packaged mesh for the
transformer, Bore~1 and Bore~2, total rank 100; first refinement for Bore~3, rank 200 &
Complete on Cartesian grids ($S=12$ at eight slabs, $S=6$ at 16 and 32) and Bore~3
($S=50$); truncated elsewhere \\
Sliced Hessian & Table~\ref{tab:v20_operation} (steady and trajectory sequences) &
Principal submatrix of the preassembled Hessian & Packaged mesh, total rank 100; first
refinement for Bore~3, rank 200 & Truncated for the transformer day, overloads, drive
cycle and four-slab shuffled sequences. The steady sequences have no temporal
factors \\
\bottomrule
\end{tabularx}
\end{table}

\begin{algorithm}[htbp]
\caption{The sliced-Hessian procedure builds one reference and reuses it for every inactive-set solve of a query sequence.}
\label{alg:final_procedure}
\begin{algorithmic}[1]
\Require Queries $q=1,\dots,N_q$ with target, bound, load and flow scale; prescribed reference mesh and rank $r$, or $S$ spatial modes with all $N_t$ temporal factors each
\State Assemble the reduced Hessian at the nominal flow
\State Form the closed-form Laplacian modes on a Cartesian grid, or compute the lowest generalized eigenvectors of the coarse reduced Hessian at the nominal flow with residuals at most $10^{-7}$ and prolong them to the target mesh
\State For a trajectory, keep the $r$ lowest spatial--temporal pairs, or all $N_t$ temporal factors of each of $S$ spatial modes
\State Store the reference on the GPU
\For{$q=1,\dots,N_q$}
  \If{the flow scale of query $q$ differs from that of the current operator}
    \State Reassemble the transport and the reduced Hessian; keep the reference
  \EndIf
  \State Form the right-hand side from the target and load of query $q$; for a chained window, take the initial state from the final state of window $q-1$
  \State Start from the previous converged state and active set; the first query starts with every constraint active
  \Repeat
    \State Extract the inactive principal submatrix of the current Hessian and restrict the reference on the GPU
    \State Solve with Algorithm~\ref{alg:revision_deflation}; verify the relative residual at most $10^{-10}$
    \State Update the multipliers and the active set for the bound of query $q$
  \Until{every KKT measure is at most $10^{-8}$}
\EndFor
\end{algorithmic}
\end{algorithm}
